\documentclass[10pt]{article}
\usepackage[a4paper,margin=2cm]{geometry}
\usepackage{lmodern,microtype,amsmath,booktabs,caption,siunitx,xcolor,hyperref}
\usepackage{pgfplots}
\pgfplotsset{compat=1.18}
\hypersetup{colorlinks,linkcolor=black,urlcolor=blue!50!black}
\captionsetup{font=small,labelfont=bf,skip=4pt}
\newcommand{\qualityWorlds}{270}
\newcommand{\optimizeWorstGap}{1.5e-12}
\newcommand{\greedyPlainMedianGap}{9.98}
\newcommand{\greedyPlainWorstGap}{89.59}
\newcommand{\threads}{16}
\newcommand{\originalTableDeviation}{4e-13}
\newcommand{\maxContenders}{4}
\newcommand{\bianchiErrorScenario}{0.005}
\newcommand{\bianchiErrorLarge}{0.018}
\newcommand{\hiddenError}{0.002}
\newcommand{\alohaError}{0.218}
\newcommand{\estimateLossError}{0.49}
\newcommand{\estimateCodecError}{0.10}
\newcommand{\isolatedCases}{48}
\newcommand{\isolatedLossError}{0.17}
\newcommand{\isolatedCodecError}{0.03}
\newcommand{\voiceExactCases}{80}
\newcommand{\voiceExactGap}{5e-13}
\newcommand{\hOneWorlds}{50}
\newcommand{\hOneWins}{49}
\newcommand{\hOneTies}{1}
\newcommand{\hOneLosses}{0}
\newcommand{\hOneMedian}{3.12}
\newcommand{\hOneP}{1.8e-15}
\newcommand{\hOneBWorlds}{50}
\newcommand{\hOneBWins}{50}
\newcommand{\hOneBTies}{0}
\newcommand{\hOneBLosses}{0}
\newcommand{\hOneBMedian}{136.40}
\newcommand{\hOneBP}{8.9e-16}
\newcommand{\hTwoWorlds}{50}
\newcommand{\hTwoWins}{50}
\newcommand{\hTwoTies}{0}
\newcommand{\hTwoLosses}{0}
\newcommand{\hTwoMedian}{27.19}
\newcommand{\hTwoP}{8.9e-16}
\newcommand{\hThreeKept}{200}
\newcommand{\hThreeTrials}{200}
\newcommand{\hTwoCodecCut}{26.1}

\newcommand{\figtable}[2]{\par\medskip\noindent\begin{minipage}{\linewidth}\centering\captionof{table}{#1}#2
  \end{minipage}\par\medskip}

\title{Channel assignment and voice QoS in synthetic 802.11 networks:\\ results}
\author{}
\date{\today}

\begin{document}
\maketitle

\noindent All numbers below are produced by \texttt{paper/results.jl} on synthetic worlds (\threads{} threads);
the model and methods are described in the companion document \texttt{model.pdf}.
A world places up to 18 access points (APs) at random, each with its clients, and an assignment picks one channel per AP.
The throughput objective is proportional fairness, $\sum_c \log r_c$ over client rates $r_c$, measured in nats.

\section{Channel assignment against the exact optimum}

Each method is scored against the exact optimum from the subset dynamic program, on \qualityWorlds{} worlds:
10, 14 and 18 APs, 4 and 16 clients per AP, ten seeds per setting for 2, 3, 4 and 6 channels at 40\,Mbps demand,
and five seeds for 3 channels at 3\,Mbps (``light'', where APs are not saturated and the airtime fixed point is active).
These worlds use an area of $100\sqrt{n/18}$\,m per side, 8\,dB shadowing and a 0.001\,Mbps rate floor; the voice
scenarios use the default model. Annealing runs 50\,000 moves.

\figtable{Share of worlds solved exactly (gap below $10^{-9}$ nats), by number of channels.}{\begin{tabular}{lrrrrr}
\toprule
Method & \multicolumn{1}{c}{2} & \multicolumn{1}{c}{3} & \multicolumn{1}{c}{4} & \multicolumn{1}{c}{6} & \multicolumn{1}{c}{3, light} \\
\midrule
Graph colouring & 2\,\% & 2\,\% & 0\,\% & 0\,\% & 10\,\% \\
Greedy, single pass (original) & 2\,\% & 0\,\% & 3\,\% & 3\,\% & 37\,\% \\
Greedy with local search & 63\,\% & 38\,\% & 30\,\% & 43\,\% & 73\,\% \\
Annealing, fixed $T_0 = 5$, no polish & 100\,\% & 95\,\% & 90\,\% & 72\,\% & 100\,\% \\
Annealing, calibrated $T_0$, polish & 100\,\% & 100\,\% & 92\,\% & 80\,\% & 97\,\% \\
\texttt{optimize\_assignment} & 100\,\% & 100\,\% & 100\,\% & 100\,\% & 100\,\% \\
\bottomrule
\end{tabular}
}

\figtable{Gap to the optimum in nats: median (worst), by number of channels.}
{\small\setlength{\tabcolsep}{4pt}\begin{tabular}{lrrrrr}
\toprule
Method & \multicolumn{1}{c}{2} & \multicolumn{1}{c}{3} & \multicolumn{1}{c}{4} & \multicolumn{1}{c}{6} & \multicolumn{1}{c}{3, light} \\
\midrule
Graph colouring & 27.56 (126.77) & 40.94 (270.56) & 67.49 (353.05) & 95.56 (437.04) & 38.17 (177.81) \\
Greedy, single pass (original) & 10.30 (89.59) & 12.88 (72.33) & 11.48 (60.76) & 6.09 (53.02) & 2.98 (40.81) \\
Greedy with local search & 0.00 (43.26) & 1.12 (21.45) & 1.20 (38.15) & 0.34 (9.59) & 0.00 (11.88) \\
Annealing, fixed $T_0 = 5$, no polish & 0.00 (0.00) & 0.00 (3.49) & 0.00 (1.55) & 0.00 (5.68) & 0.00 (0.00) \\
Annealing, calibrated $T_0$, polish & 0.00 (0.00) & 0.00 (0.00) & 0.00 (3.82) & 0.00 (3.09) & 0.00 (0.13) \\
\texttt{optimize\_assignment} & 0.00 (0.00) & 0.00 (0.00) & 0.00 (0.00) & 0.00 (0.00) & 0.00 (0.00) \\
\bottomrule
\end{tabular}
}

\noindent \texttt{optimize\_assignment} reaches the optimum in every world (worst gap \num{\optimizeWorstGap} nats).
The original single-pass greedy is \greedyPlainMedianGap{} nats short in the median and up to \greedyPlainWorstGap{} nats.

\section{Speed}

\figtable{Median wall time in ms, 16 clients per AP, 3 channels. ``Original table'' is the original single-threaded
subset-table algorithm (same output to within \num{\originalTableDeviation}); it does not apply to light demand, which needs
the per-group airtime fixed point. ``Voice-aware'' is the optimizer with the voice objective (64 calls per AP, 75\,\%
background load).}{\begin{tabular}{lrrrrrr}
\toprule
Demand & APs & Original table & New table & Optimize & Exact & Voice-aware \\
\midrule
heavy & 12 & 18.3 & 6.3 & 7.6 & 6.3 & 11.1 \\
heavy & 14 & 93.3 & 6.1 & 14.0 & 11.4 & 35.1 \\
heavy & 16 & 401.2 & 29.4 & 24.8 & 24.2 & 130.7 \\
heavy & 18 & 1799.7 & 73.8 & 72.4 & 97.7 & 564.4 \\
light & 12 & -- & 11.2 & 13.0 & 11.4 & 25.4 \\
light & 14 & -- & 48.4 & 52.1 & 49.5 & 108.9 \\
light & 16 & -- & 193.0 & 192.5 & 196.0 & 412.6 \\
light & 18 & -- & 1769.7 & 1165.7 & 1449.1 & 2127.4 \\
\bottomrule
\end{tabular}
}

\section{Model validation}

\begin{center}
\begin{tikzpicture}
\begin{axis}[width=0.46\linewidth,height=5cm,xlabel={contenders $n$},ylabel={collision probability},
             legend pos=south east,legend style={font=\small},title={Bianchi contention}]
\addplot[thick,blue] table[x=n,y=model]{generated/bianchi.dat};
\addplot[only marks,mark=o,red] table[x=n,y=simulation]{generated/bianchi.dat};
\legend{model,slot simulation}
\end{axis}
\end{tikzpicture}\hfill
\begin{tikzpicture}
\begin{axis}[width=0.46\linewidth,height=5cm,xlabel={hidden AP airtime $u$},ylabel={overlap probability},
             legend pos=south east,legend style={font=\small},title={Hidden transmitter}]
\addplot[thick,blue] table[x=u,y=model]{generated/hidden.dat};
\addplot[only marks,mark=o,red] table[x=u,y=simulation]{generated/hidden.dat};
\addplot[dashed,gray] table[x=u,y=aloha]{generated/hidden.dat};
\legend{model,simulation,ALOHA $1-e^{-2u}$}
\end{axis}
\end{tikzpicture}
\captionof{figure}{Collision models against simulation. The voice scenarios reach at most \maxContenders{} contenders (the AP
plus the co-channel APs it senses); up to that, Bianchi's fixed point is within \bianchiErrorScenario{} of a slot-level DCF
simulation, and within \bianchiErrorLarge{} up to 10. The single-transmitter overlap formula is within \hiddenError{};
the ALOHA formula it replaced was off by up to \alohaError{}.}
\end{center}

\figtable{Steady-state voice estimate (used by the voice-aware objective) against \texttt{simulate\_voice}: worst
absolute error in loss (percentage points) / codec rate (kbps) over 20 assignments per cell.}{\begin{tabular}{lrrrr}
\toprule
Scenario & \texttt{:fifo} & \texttt{:deadline} & \texttt{:adaptive} & \texttt{:combined} \\
\midrule
Branch office & 0.00 / 0.00 & 0.00 / 0.00 & 0.00 / 0.00 & 0.00 / 0.00 \\
Open-plan office & 0.22 / 0.00 & 0.00 / 0.00 & 0.09 / 0.02 & 0.00 / 0.00 \\
High-density call centre & 0.49 / 0.00 & 0.15 / 0.00 & 0.32 / 0.07 & 0.00 / 0.01 \\
Emergency overload stress & 0.39 / 0.00 & 0.34 / 0.00 & 0.45 / 0.10 & 0.06 / 0.06 \\
\bottomrule
\end{tabular}
}

\noindent The queue simulator matches the closed-form steady state of an isolated AP in all \isolatedCases{} cases
(loss within \isolatedLossError{} points, codec within \isolatedCodecError{} kbps), and the voice-aware optimizer matches an
exhaustive search of its objective in all \voiceExactCases{} cases (worst gap \num{\voiceExactGap}).

\section{Voice}

\figtable{Voice outcomes on throughput-optimal assignments, 20 worlds per scenario, 10\,s simulated.}
{\small\begin{tabular}{llrrrr}
\toprule
Scenario & Policy & Loss (median) & Loss (worst) & p95 delay (ms) & Codec (kbps) \\
\midrule
Branch office & \texttt{:fifo} & 0.01\,\% & 0.65\,\% & 7.8 & 64.0 \\
 & \texttt{:deadline} & 0.01\,\% & 0.65\,\% & 0.4 & 64.0 \\
 & \texttt{:adaptive} & 0.01\,\% & 0.65\,\% & 7.8 & 64.0 \\
 & \texttt{:combined} & 0.01\,\% & 0.65\,\% & 0.4 & 64.0 \\
Open-plan office & \texttt{:fifo} & 0.16\,\% & 2.10\,\% & 17.5 & 64.0 \\
 & \texttt{:deadline} & 0.16\,\% & 1.03\,\% & 1.8 & 64.0 \\
 & \texttt{:adaptive} & 0.16\,\% & 1.44\,\% & 17.5 & 64.0 \\
 & \texttt{:combined} & 0.16\,\% & 1.03\,\% & 1.8 & 64.0 \\
High-density call centre & \texttt{:fifo} & 11.73\,\% & 19.31\,\% & 140.0 & 64.0 \\
 & \texttt{:deadline} & 1.19\,\% & 3.45\,\% & 8.5 & 64.0 \\
 & \texttt{:adaptive} & 5.05\,\% & 9.83\,\% & 140.0 & 36.5 \\
 & \texttt{:combined} & 1.19\,\% & 3.45\,\% & 8.5 & 64.0 \\
Emergency overload stress & \texttt{:fifo} & 40.85\,\% & 46.89\,\% & 140.0 & 64.0 \\
 & \texttt{:deadline} & 5.55\,\% & 15.09\,\% & 140.0 & 64.0 \\
 & \texttt{:adaptive} & 20.85\,\% & 27.84\,\% & 140.0 & 12.2 \\
 & \texttt{:combined} & 0.53\,\% & 2.90\,\% & 37.6 & 59.6 \\
\bottomrule
\end{tabular}
}

\figtable{Throughput-only against voice-aware assignment (same policy for both): median loss, worlds where voice goodput
rises / falls, and the cost to data clients in nats per client.}{\small\setlength{\tabcolsep}{4pt}\begin{tabular}{llrrrrr}
\toprule
Scenario & Policy & Loss before & Loss after & Goodput up/down & Cost (median) & Cost (max) \\
\midrule
Branch office & \texttt{:fifo} & 0.01\,\% & 0.01\,\% & 1 / 0 & 0.000 & 0.000 \\
 & \texttt{:deadline} & 0.01\,\% & 0.01\,\% & 1 / 0 & 0.000 & 0.000 \\
 & \texttt{:adaptive} & 0.01\,\% & 0.01\,\% & 1 / 0 & 0.000 & 0.000 \\
 & \texttt{:combined} & 0.01\,\% & 0.01\,\% & 1 / 0 & 0.000 & 0.000 \\
Open-plan office & \texttt{:fifo} & 0.16\,\% & 0.08\,\% & 2 / 0 & 0.000 & 0.034 \\
 & \texttt{:deadline} & 0.16\,\% & 0.16\,\% & 0 / 0 & 0.000 & 0.000 \\
 & \texttt{:adaptive} & 0.16\,\% & 0.06\,\% & 3 / 0 & 0.000 & 0.034 \\
 & \texttt{:combined} & 0.16\,\% & 0.16\,\% & 0 / 0 & 0.000 & 0.000 \\
High-density call centre & \texttt{:fifo} & 11.73\,\% & 8.42\,\% & 17 / 0 & 0.039 & 0.092 \\
 & \texttt{:deadline} & 1.19\,\% & 0.58\,\% & 9 / 0 & 0.000 & 0.033 \\
 & \texttt{:adaptive} & 5.05\,\% & 2.17\,\% & 19 / 0 & 0.084 & 0.181 \\
 & \texttt{:combined} & 1.19\,\% & 0.58\,\% & 9 / 0 & 0.000 & 0.033 \\
Emergency overload stress & \texttt{:fifo} & 40.85\,\% & 40.36\,\% & 16 / 0 & 0.010 & 0.042 \\
 & \texttt{:deadline} & 5.55\,\% & 3.38\,\% & 19 / 0 & 0.034 & 0.078 \\
 & \texttt{:adaptive} & 20.85\,\% & 19.79\,\% & 16 / 0 & 0.023 & 0.076 \\
 & \texttt{:combined} & 0.53\,\% & 0.65\,\% & 19 / 0 & 0.032 & 0.088 \\
\bottomrule
\end{tabular}
}

\noindent Paired sign tests over \hOneWorlds{} worlds: deadline scheduling beats adaptation alone on loss in \hOneWins{}
worlds (\hOneTies{} tied, \hOneLosses{} lost; median \hOneMedian{} points, $p = \num{\hOneP}$) and cuts p95 delay in
\hOneBWins{} (median \hOneBMedian{}\,ms, $p = \num{\hOneBP}$). Under voice overload, priority with codec adaptation beats
fixed-rate priority in \hTwoWins{} of \hTwoWorlds{} worlds (median \hTwoMedian{} points, $p = \num{\hTwoP}$) at a median
codec reduction of \hTwoCodecCut{}\,kbps, and keeps at least 99\,\% of the codec rate in \hThreeKept{} of \hThreeTrials{}
normal-load trials.

\section{Limits}
Worlds, traffic and shadowing are synthetic and static. The voice model is an aggregate work queue, not a packet-level
802.11 simulator: Bianchi treats partly busy APs as fractional contenders, and each AP's own clients count as one contender,
which likely understates collisions in call-heavy cells. The exact methods cover up to 18 APs. External validation needs
ns-3, packet traces or hardware.

\end{document}
